\documentclass[11pt]{article}
\usepackage[T1]{fontenc}
\usepackage{lmodern,microtype}
\usepackage{amsmath,amssymb}
\usepackage[letterpaper,margin=1in]{geometry}
\usepackage{booktabs,array,enumitem}
\usepackage[colorlinks=true,linkcolor=blue!50!black,urlcolor=blue!50!black]{hyperref}
\newcommand{\GL}{\mathrm{GL}}
\newcommand{\A}{\mathbb A}
\title{Referee report on \texttt{mjl\_part4\_draft3.tex}\\
\normalsize ``A metaplectic Jacquet--Langlands transfer for division
algebras of prime degree, IV'', Draft~3}
\author{Claude (Anthropic), for David Joyner}
\date{November 7, 2026}
\begin{document}
\maketitle

\section*{0. Scope}
\paragraph{Read in full.} Draft~3 (1090 lines; 16 pages; it compiles
cleanly) and its revision log.

\paragraph{Compiled numbering.}
\begin{itemize}
\item Proposition~2.1;
\item Theorem~3.1, Example~3.2, Remark~3.3;
\item Lemma~4.1, Proposition~4.2, Remark~4.3;
\item Proposition~5.1, Theorem~5.2;
\item Proposition~6.1;
\item Proposition~7.1;
\item Definition~8.1, Lemma~8.2, Proposition~8.3.
\end{itemize}

\paragraph{Which versions I could check.} Draft~3 now cites Part~I,
Draft~7, Part~II, Draft~7, and Part~III, Draft~9. I have only Draft~6
of each.
\begin{itemize}
\item Every cited number agrees with Draft~6:
\begin{itemize}
\item Part~I: Propositions~2.3, 2.6, 3.1, 3.5, 5.5 and~6.1; Lemma~5.1;
Example~3.7; Theorem~6.6;
\item Part~II: Definition~2.1, Theorem~6.2;
\item Part~III: Lemma~3.3; Theorems~4.3, 7.6 and~9.1; Corollaries~5.1
and~5.2; Lemma~7.5; Proposition~11.3.
\end{itemize}
\item I could not check Part~III, Corollary~7.8 (new in Draft~9), or the
new subtitles of Parts~I and~II.
\item \S8.4 now cites ``Part~I, Proposition~2.3 and its archimedean
remark''. I could not identify that remark in Draft~6. The
complex-place splitting is (R4) of Part~I, Appendix~B, according to the
revision log. I did not read that appendix.
\end{itemize}

\paragraph{External sources.} The only new external citation is
Clozel--Delorme. Its bibliographic data agree with the standard
citation as far as I know; I did not check the paper. I have not checked
Li's paper, so I cannot confirm that its Paley--Wiener hypothesis is
stated only at \emph{real} places.

\section*{Summary verdict}
Every point of my Draft~2 report has been addressed correctly. I
rechecked each repair:
\begin{itemize}
\item the new Haar-modulus proof (Proposition~4.2), which is now correct
and complete;
\item the justification of $u(x)=1$;
\item Remark~3.3 and its field-torus proof;
\item the central-character discussion in Proposition~7.1;
\item the Li/Clozel--Delorme paragraph;
\item the merged spectral remark.
\end{itemize}

\textbf{Sections~2--7 are now correct and essentially ready.} I have only
minor editorial points on them (\S3).

\textbf{\S8 is the new part, and it is correct as stated.} I found no
error in Definition~8.1, Lemma~8.2 or Proposition~8.3. But I have four
substantive comments (\S1). One affects how Proposition~8.3 should be
read: its product condition~(iii) is a \emph{normalization}, not an
existence problem. Another identifies a case that looks provable now:
$r=2$.

\paragraph{On ``getting closer to the goals''.} Two kinds of progress
should be kept apart.
\begin{itemize}
\item \emph{As a paper,} Draft~3 is close to finished. Its claims are
correct, correctly labeled and correctly attributed.
\item \emph{As progress on the series' two open goals,} it has not
proved either of them:
\begin{itemize}
\item integrality of $L_D$ at wild places;
\item the comparison of $D^\times$ with its cover.
\end{itemize}
What it has done, and this is real progress, is reduce the second goal
to a short list of precise local statements at split places. For $r=2$
part of that list already appears to hold (\S1.3).
\end{itemize}

\section*{1. Section 8}

\subsection*{1.1 The product condition in Proposition~8.3 is a normalization}
Suppose (i) and (ii) hold, with every $\lambda_v$ independent of the
class. Then the proof of Proposition~8.3 gives, with no further
hypothesis,
\[
I_{\widetilde H}(\widetilde f)=\Lambda\,I_H(f),\qquad
\Lambda=\prod_{v\notin S}\lambda_v\prod_{v\in S}b_{D,v}.
\]
Hypothesis (iii), $\Lambda=1$, can always be arranged by rescaling one
test function. So it is not a constraint on existence.

What matters is the \emph{value} of $\Lambda$ under canonical
normalizations: $L_D$ (unitary) at $S$, and the Flicker--Kazhdan unit
elements and matching at $v\notin S$. Integrality of the global
multiplicities then constrains the local character coefficients.

I recommend restating the conclusion as $I_{\widetilde H}=\Lambda I_H$,
with $\Lambda$ written out, and noting that $\Lambda$ is computable
once the split-place normalization is fixed. The real hypothesis in (i)
is \emph{uniformity}: $\lambda_v$ must not depend on $s_v$, and in
particular not on the torus type $k_v$.

\subsection*{1.2 The trace formulas are simpler than Proposition~8.3 assumes}
Both sides are compact-quotient trace formulas: $D$ is a division
algebra, and the cover is compact modulo the central preimage. Since
$r$ is prime, every noncentral element of $D^\times(K)$ is regular
elliptic. Consequently:
\begin{itemize}
\item there are no singular noncentral classes;
\item there is no continuous spectrum;
\item absolute convergence holds for every smooth compactly supported
test function;
\item the only term the support condition must remove is the identity
term.
\end{itemize}
The convergence clause in (iii) and the talk of ``singular terms'' can
therefore be dropped.

Better: \emph{keep} the identity term, and use it as an independent
check on $\Lambda$. The relevant facts:
\begin{itemize}
\item the adelic quotients $Z(\A)\backslash\overline{G_D}(\A)$ modulo
$\iota(D^\times(K))$ and $\A^\times D^\times(K)\backslash
D^\times(\A)$ are homeomorphic, so they have equal volumes;
\item at $v\in S$, $\widetilde f_v(1)=c_{D,v}f_v(1)$ (Part~III,
Corollary~9.3).
\end{itemize}
On the regular terms the constant at $v\in S$ is $b_{D,v}$, while on
the identity term it is $c_{D,v}=b_{D,v}|n|_v^{-r(r-1)/2}$. So the
identity terms must supply the extra factor
$\prod_{v\in S}|n|_v^{-r(r-1)/2}$, balanced by the split places.

The ratio $\Delta(t)/\Delta(t^n)$ tends to $|n|_v^{-r(r-1)/2}$ near
scalars for \emph{every} torus type. The proof of Part~III, Lemma~3.3,
only uses $(1-u^n)/(1-u)\to n$. This suggests that the split-place
identity values contribute $|n|_v^{-r(r-1)/2}$ as well. If so, the
product formula $\prod_v|n|_v=1$ balances the identity terms, provided
the complex places are included, where $|n|_v=n^2$.

This is heuristic. But it is a concrete consistency test that the
draft could record, and it uses the identity-term comparison of
Part~I, Proposition~6.1.

\subsection*{1.3 Definition~8.1 names the wrong ``in particular''}
The example given, $J(\xi x,f_v)=J(x,f_v)$ for scalar $\xi\in\mu_n$,
is \emph{automatic}. The ordinary test function has central character
$\theta^{-1}$, and $\theta(\xi)=\tilde\theta(\xi^*)=1$ (\S2.2). So
Lemma~8.2's ``in particular'' case has no content.

The real constraint is $x'=\zeta x$ with $\zeta\in\mu_n(E_v)$
\emph{nonscalar}, i.e.\ $\zeta$ a nontrivial element of
$\ker P_n\cong\mu_n^{k_v}/\mu_n$ (Proposition~4.2). Then
$(x')^*=x^*$. One checks this by noting $\zeta^*=1$: every symbol
$(\zeta_a,\zeta_b)_n$ is a power of $(\zeta_n,\zeta_n)_n=1$. The
condition is then $J(\zeta x,f_v)=J(x,f_v)$. Please state this case
explicitly. It exists only for non-field tori ($k_v\ge2$).

\paragraph{The unit element satisfies it on split tori.} Let $t$ be
regular in $T(\mathcal O_v)$ for the split torus, with $v\nmid n$, and
use the standard measures. Then
\[
O(t,1_{K_v})=\prod_{\alpha>0}|1-\alpha(t)|^{-1}
\qquad\text{and}\qquad
\Delta(t)=\prod_{\alpha>0}|1-\alpha(t)|,
\]
so the normalized orbital integral is $J(t,1_{K_v})=1$ for all such
$t$. In particular $J(\zeta t)=J(t)$. Hence:
\begin{itemize}
\item \emph{For $r=2$} the only tori are split ($k_v=2$) and fields
($k_v=1$, where the condition is vacuous). So the unramified unit
elements are FK-compatible at every unramified split place. This makes
$r=2$, $v\nmid n$, a realistic first case in which to prove hypothesis
(i) of Proposition~8.3 outright.
\item \emph{For $r\ge3$} there are mixed tori $F\times E'$, and
compatibility of $1_{K_vZ_v}$ there is not clear. Translating by a
nonscalar root of unity moves an element near the center away from it,
and the normalized orbital integral of the unit element is not constant
near the center. This should be checked, or extracted from
[FK, \S\S11--12], before (i) is asserted with $\lambda_v=1$ at almost
all places.
\end{itemize}

\subsection*{1.4 Two small gaps in the proof of Proposition~8.3}
\begin{enumerate}[label=(\alph*)]
\item \emph{Rational lift versus star.} The proof uses that the
covered term at the rational class $\iota(s^n)$ is the local star
$s_v^*$ at every $v$. This holds because the rational splitting $\iota$
is a homomorphism: $\iota(s^n)=\iota(s)^n$, and $\iota(s)_v$ is a lift
of $s_v$. Add this sentence.
\item \emph{Complex places.} Hypothesis~(i) is imposed at the complex
places, where the cover splits. There, the compatibility and the
pointwise identity say that $J(x,f_\infty)$ is a pullback of an
orbital-integral function by $x\mapsto x^n$. So the same
singular-power issue as in the ``regularity and germ obligation''
paragraph arises at $x=\mathrm{diag}(1,\zeta,\ldots)$. That paragraph
should say it applies at the archimedean places too.
\end{enumerate}

\subsection*{1.5 Lemma~8.2}
Lemma~8.2 is correct, and its proof is a one-line consequence of the
central transformation law. One hypothesis is redundant: ``The
corresponding identity holds also at central translates'' follows if
the region is required to be stable under $Z_v$, which it should be
anyway. With \S1.3, the lemma's content is the nonscalar kernel case.

\section*{2. Sections 2--7: rechecks}
\begin{itemize}
\item \emph{Proposition~4.2.} The new proof is correct.
\begin{itemize}
\item $P_n^{-1}(U^n)=K_0U$ and $\mathrm{vol}(K_0U)=|K_0|\,\mathrm{vol}(U)$.
\item $[U:U^n]=|n|_F^{-(r-1)}$ from the $\mathbb Z_p$-rank
$(r-1)[F:\mathbb Q_p]$.
\end{itemize}
One point deserves a clause: the image of principal units in
$E^\times/F^\times$ can contain torsion, e.g.\ $\zeta_p\in E\setminus F$.
``Sufficiently small'' handles this, because with compatible levels
$U_E^{(em)}\cap F=U_F^{(m)}$ the quotient lattice is saturated, hence
torsion-free. Remark~4.3 and its $\mathbb Z$ counterexample are now
right.
\item \emph{\S2.2.} The argument $u(x)\in\{\pm1\}\cap\mu_n=\{1\}$ is
correct. The global statement $\theta|_{K^\times}=1$ is correct.
\item \emph{Remark~3.3.} Correct, including the reason
$\zeta_3\notin F^{\times3}$ (no $\zeta_9$ in $\mathbb Q_3(\zeta_3)$),
and the field-torus proof.
\item \emph{Proposition~7.1.} Correct. One suggestion: the compatibility
$\omega_\pi(a)=\omega_{\widetilde\pi}(a^*)$ need not be assumed.
Replace $x$ by $ax$ in Zou's (3.2), with $a$ central, using
$(ax)^*=a^*x^*$ and that the roots of $(ax)^n$ are $at$. This gives
the compatibility whenever $\theta_{\widetilde\pi}$ is not identically
zero on the elements $x^*$ with $x^n$ regular. One line would make the
proposition self-contained.
\item \emph{\S8.4.} The argument is now internally consistent: no real
places, the covers split at complex places, and Clozel--Delorme covers
the linear case. See \S0 for what I could not verify.
\end{itemize}

\section*{3. Editorial points}
\begin{enumerate}[label=(\arabic*)]
\item \emph{Undefined symbol.} $I_F$ in the display of wild constants in \S6.3 (source line 616)
is never defined. Write
``$[E^\times:F^\times E^{\times3}]=9$'' or drop it.
\item \emph{``New''.} The title of \S4, ``A \emph{new} local power-map
calculation'', overstates. The index formula is classical; what is
useful here is the Haar modulus and the bookkeeping. Drop ``new''.
\item \emph{Duplicated display.} The display for $c(\gamma)$ appears
twice, in Remark~5.3 and \S8.1. Keep one.
\item \emph{Labels.} The status convention now matches Part~I. Its claim
``as in Parts~I--III'' depends on Part~III, Draft~9, which I have not
seen. Part~III, Draft~6, used \textbf{[O]} for ``open''. Please
confirm.
\item \emph{Wording.} The acknowledgments call the paper a ``research
note''; elsewhere it is ``Part~IV''. Otherwise the acknowledgments are
fine, and the drafting history is now out of the body.
\item \emph{Abstract.} It is accurate. It could add a sentence describing
Proposition~8.3 as reducing the $D^\times$-versus-cover comparison to
explicit split-place conditions.
\end{enumerate}

\section*{4. Recommendation}
\begin{itemize}
\item \emph{Sections~2--7:} accept, after the editorial points.
\item \emph{Section~8:} minor revision:
\begin{itemize}
\item restate Proposition~8.3 with the constant $\Lambda$ (\S1.1);
\item simplify its hypotheses using compactness (\S1.2);
\item correct the ``in particular'' of Definition~8.1 and Lemma~8.2
(\S1.3);
\item add the two sentences of \S1.4.
\end{itemize}
\end{itemize}

\paragraph{The most promising next mathematical steps.}
\begin{enumerate}[label=(\alph*)]
\item \emph{Prove (i) of Proposition~8.3 for $r=2$ at places
$v\nmid n$.} The unit element is FK-compatible there (\S1.3), and only
two torus types occur. A pointwise identity
$J(x^*,\widetilde f_v)=\lambda_vJ(x,f_v)$ for the unit elements is then
a fundamental-lemma statement. I expect [FK, \S12] to supply it, since
$(n,N_2)=(n,1)=1$ holds automatically, but I have not checked that
FK's statement there has exactly this pointwise form.
\item \emph{Compute $\Lambda$ in that case,} and test it against the
identity-term heuristic of \S1.2.
\item \emph{Carry out the finite computation for
$(2,3,\mathbb Q_3(\zeta_3))$} (\S6.3). It is still the cleanest way to
decide the first wild case.
\end{enumerate}
Items (a) and (b) would give the first unconditional instance of the
$D^\times$-versus-cover comparison, with test functions restricted at
$S$ and at the places dividing~$n$.

\end{document}
